% Sentinel -- arXiv preprint draft.
% Every number in this file must carry an evidence ID from paper/evidence.md,
% written as \ev{B1}. Do not quote a number that has no row there.
\documentclass[11pt]{article}

\usepackage[margin=1in]{geometry}
\usepackage[T1]{fontenc}
\usepackage{lmodern}
\usepackage{microtype}
\usepackage{amsmath,amssymb}
\usepackage{booktabs}
\usepackage{graphicx}
\usepackage{xcolor}
\usepackage[numbers,sort&compress]{natbib}
\usepackage[hidelinks]{hyperref}

\newcommand{\todo}[1]{\textcolor{red}{[TODO: #1]}}
\newcommand{\ev}[1]{\textcolor{gray}{\scriptsize[#1]}}   % evidence pointer; strip before submission

\title{Sentinel: Latent-Fusion Agents for Early Warning of ICU Deterioration, Validated on 208 External Hospitals}
\author{\todo{author names and affiliations} \\ PES University}
\date{\todo{date}}

\begin{document}
\maketitle

\begin{abstract}
\todo{Write last. Problem; approach (agents pass latents + calibrated risks to two
attention--TCN heads); main result \ev{B1,B2,B4}; operating-point result \ev{C1,C4};
external validation \ev{G1,G2}; deployment footprint \ev{I1,I2}; one-line limits.}
\end{abstract}

% ---------------------------------------------------------------------------
\section{Introduction}\label{sec:intro}
% - Clinical problem: deterioration in the ICU, alarm fatigue, NEWS2 as the
%   bedside standard \todo{cite RCP 2017}.
% - ML early-warning scores exist (eCART, YAIB, HiRID benchmarks); gradient-
%   boosted trees on tabular summaries are the bar \cite{churpek2025ecartv5,vandewater2024yaib,yeche2021hirid}.
% - Multi-agent LLM systems are popular but pass text between agents; for
%   numeric physiology this is lossy \ev{D1}. CLAIM WORDING DEPENDS ON D2 RE-RUN.
% - Contributions (bullets): architecture; evaluation protocol (frozen
%   thresholds, matched alarm budgets, patient-clustered CIs, 10 seeds);
%   external validation; deployment + explanation measurements; negative results.
\todo{draft}

% ---------------------------------------------------------------------------
\section{Related work}\label{sec:related}
% 2.1 Clinical early-warning scores (NEWS2, eCART) and ML EWS on MIMIC/eICU.
% 2.2 ICU benchmarks: YAIB, HiRID-ICU-Benchmark, MIMIC-Extract
%     \cite{vandewater2024yaib,yeche2021hirid,wang2020mimicextract}.
% 2.3 Same-label prior work: Sadanandan 2026 \cite{sadanandan2026multimodal}.
% 2.4 Temporal convolutional networks; additive attention; mixture of experts.
% 2.5 Multi-agent LLM systems in medicine \todo{verify and add citations}.
\todo{draft}

% ---------------------------------------------------------------------------
\section{Problem setup}\label{sec:setup}

\subsection{Cohort and outcome}
We use MIMIC-IV v3.1 \cite{johnson2023mimiciv}, accessed through PhysioNet
under its credentialed data use agreement. We include every ICU stay of at
least 12 hours and represent it on an hourly grid that starts at ICU
admission and is capped at 14 days. At each hour $t$ the target is a
composite outcome: whether the first of three events occurs in the following
24 hours \ev{A2}:
\begin{enumerate}
  \item onset of invasive mechanical ventilation (including tracheostomy),
  \item onset of a vasoactive agent (norepinephrine, epinephrine, dopamine,
        phenylephrine, vasopressin, dobutamine or milrinone), or
  \item in-hospital death.
\end{enumerate}
A stay is censored at its first event; later hours are not prediction
opportunities. Transfer to the ICU, the usual ward outcome, is undefined for
patients who are already there.

Without further restriction, 73\% of first events fall within six hours of
ICU admission \ev{A3}. These are mostly patients who arrive already ventilated
(for example, from surgery) or who start a vasopressor on arrival: the event
is the reason for admission, not a deterioration that could have been
anticipated. We therefore require a stay to be event-free through hour~6 and
begin scoring there. Measurements from hours 0--5 enter only as
carried-forward context at hour~6. The resulting cohort has 56{,}157 stays and
2{,}460{,}983 scored patient-hours \ev{A1}.

\subsection{Splits}
Patients, not stays, are split 70/15/15 into training, validation and test
sets with a fixed seed, so no patient contributes hours to more than one set
\ev{A4}. The test set holds 361{,}229 patient-hours from 8{,}375 stays and
6{,}335 patients; 4.74\% of its hours are positive \ev{A1}. Validation
patients are further split in half, again by patient. The first half
($V_{\text{fit}}$) fits the fusion heads' calibration; the second
($V_{\text{freeze}}$) sets alarm thresholds and is used for any choice between
model variants. Using different patients for the two steps keeps a threshold
from being tuned on scores that were fitted to the same patients.
Individual agents use the full validation set for early stopping and for their
own calibration. The test set is scored once per reported configuration.

\subsection{Leakage control}
Every feature at hour $t$ uses only data recorded strictly before $t$
\ev{A5}. Each channel is carried forward within a stay and then shifted by one
hour. Carry-forward is bounded per channel, at roughly twice that channel's
90th-percentile charting gap: 12 hours for vitals and the Glasgow Coma Scale
(GCS), 8 hours for oxygen device and flow, and 24 hours for FiO$_2$ and
laboratory values. Older values are treated as missing. Laboratory results
are timed by when they became available in the record, not when blood was
drawn (a median of 70 minutes later). An ECG is used only if it was recorded
strictly before $t$ and is at most 72 hours old. Within an hour, vitals are
averaged and GCS takes its worst value. All normalisation statistics come from
the training set. Unit tests rewrite future inputs and check that earlier
outputs do not change.

\subsection{Metrics}
The primary metric is the area under the precision--recall curve (AUPRC),
computed as average precision,
\begin{equation}
  \mathrm{AP} = \sum_{n} (R_n - R_{n-1})\, P_n ,
\end{equation}
where $P_n$ and $R_n$ are precision and recall at the $n$-th distinct score
threshold \ev{A6}. A random ranking scores the positive rate (0.047 on the
test set). AUROC is reported alongside it for comparison with the
early-warning literature. Confidence intervals are 95\% nonparametric
bootstrap intervals that resample patients with replacement (1{,}000
resamples), so hours from the same patient stay together. Paired differences
between models are computed on the same resamples.

\paragraph{Alarms and alarm budgets.}
A model raises an alarm at hour $t$ when its risk reaches a threshold. To
compare models at equal cost, each is given the same \emph{alarm budget}: the
share of patient-hours on which it may alarm. Thresholds are set on
$V_{\text{freeze}}$ and applied unchanged to the test set. For the head-to-head
comparisons we anchor the budget to NEWS2. Its threshold is the
$(1-b)$-quantile of NEWS2 on $V_{\text{freeze}}$. Because NEWS2 is
integer-valued, its achieved volume $b'$ exceeds the nominal $b$: a 2\%
budget gives a threshold of 11 and $b' = 2.62\%$. Every other model's threshold
is then the $(1-b')$-quantile of its own $V_{\text{freeze}}$ scores. Realized
alarm rates on the test set differ slightly between models, and we report them
\ev{C2}. In the budget sweep (\S\ref{sec:results}), each model's threshold is
simply its own $(1-b)$-quantile on $V_{\text{freeze}}$.

\paragraph{Patient-level outcomes.}
A deteriorating stay is \emph{caught} if the model's first alarm falls at or
before the event, and caught \emph{early} if it falls at least 6 hours before.
Warning time is the gap between the first alarm and the event. Lead over
NEWS2 is computed on stays that both models catch, as NEWS2's first-alarm hour
minus the model's; positive values mean the model alarmed first.

% ---------------------------------------------------------------------------
\section{Sentinel}\label{sec:method}

\subsection{Agents}
Sentinel has eight agents, one per kind of data (Table~\ref{tab:agents}). Each
agent is trained on its own and, every hour, emits two things: a latent vector
$z_{a,t}$ summarising what it sees, and a calibrated risk $r_{a,t}$ (its raw
score passed through isotonic regression \todo{cite Zadrozny \& Elkan} fitted
on validation data). When an agent's data are missing it emits a learned
``no data'' value rather than an imputed reading. Because every agent shares
this output format, an agent can be added or removed and only the fusion
heads need to be retrained.

Six agents are neural networks. Vital signs are the one input that arrives as
a dense hourly sequence, so the vitals agent is a causal temporal
convolutional network (TCN) \todo{cite Bai 2018; van den Oord 2016} whose
penultimate layer is its latent. The other five read inputs that already
summarise time (the latest value and its age, or counts over 1-, 6- and
24-hour windows) and change slowly, so each is a two-layer multilayer
perceptron whose hidden layer is its latent. The remaining two agents are not
neural. The GBM agent is an XGBoost model \todo{cite Chen \& Guestrin} on 145
features from vitals and laboratory values, including 6-, 12- and 24-hour
changes and trailing minima, maxima and means. Trees have no embedding layer,
so its latent is a single number, the log-odds of its prediction. Its
training-set scores are produced out-of-fold (two patient-level folds) so that
the fusion does not learn to trust over-confident in-sample predictions. The
NEWS2 agent computes the score exactly to the Royal College of Physicians
bands (SpO$_2$ Scale~1) \todo{cite RCP 2017}, and its latent is the vector of
seven sub-scores. A missing vital scores zero, and consciousness is assessed
from GCS (alert if GCS is 15; for intubated patients, eyes~4 and motor~6).

\begin{table}[t]
\centering
\small
\caption{The eight agents. Parameter counts are trained weights
\ev{I1}.}\label{tab:agents}
\begin{tabular}{@{}l p{0.52\linewidth} l r@{}}
\toprule
Agent & Inputs & Model & Params \\
\midrule
Vitals & HR, RR, SpO$_2$, systolic BP, temperature and the age of each; GCS;
  2 flags; 6 missingness masks (19) & causal TCN & 19{,}681 \\
ECG & cart machine read: 12 intervals and axes, statement count, 12 validity
  flags, 26 statement flags; ECG age (52) & MLP & 10{,}996 \\
Cadence & counts of lab draws, charting, orders by type and cultures over 1, 6,
  24\,h; 3 acceleration ratios; 2 recency features (32) & MLP & 4{,}225 \\
Labs & 24 analytes, each with value, age and an ever-measured flag (72) & MLP
  & 3{,}713 \\
Fluid & urine output over 1--24\,h, oliguria flags, intake, net and cumulative
  balance (13) & MLP & 3{,}009 \\
Resp & oxygen-support level (0--5), airway flag, recent maximum and
  escalation, O$_2$ flow and FiO$_2$ with 6\,h changes (12) & MLP & 2{,}945 \\
GBM & 145 vitals and laboratory features with trends & XGBoost & trees \\
NEWS2 & seven sub-scores & fixed rule & 0 \\
\midrule
Fusion & both heads (\S\ref{sec:heads}) & attention + TCN & 52{,}036 \\
\bottomrule
\end{tabular}
\end{table}

\paragraph{Agent training.}
The vitals TCN uses 32 channels and three residual blocks (kernel 3,
dilations 1, 2 and 4), each with two causal convolutions, per-timestep layer
normalisation, ReLU and dropout 0.1. The MLP agents have 64 hidden units
(128 for ECG), a 32-dimensional latent and dropout 0.2. All neural agents are
trained with binary cross-entropy and Adam (learning rate $10^{-3}$), with
early stopping on validation AUPRC. Continuous inputs are standardised with
training statistics. The ECG agent sends every hour without a recent ECG the
same learned vector, so it outputs one constant risk on those hours and cannot
score well merely by noticing who had an ECG ordered. Its early stopping uses
only validation hours that have an ECG.

\subsection{Fusion heads}\label{sec:heads}
Each fusion head reads the agent bank, the eight latents and eight risks at
every hour (208 numbers), and returns one risk. Each agent's latent is first
projected to a common width $d = 32$,
\begin{equation}
  h_{a,t} = W_a z_{a,t} + b_a .
\end{equation}
Additive attention \todo{cite Bahdanau 2015} scores each agent on
its own and normalises across agents,
\begin{equation}
  e_{a,t} = v^{\top} \tanh(h_{a,t}) + c, \qquad
  \alpha_{a,t} = \frac{\exp(e_{a,t})}{\sum_{a'} \exp(e_{a',t})} .
\end{equation}
The weights are recomputed every hour for every patient. Agents are never
compared with one another, so this is not a Transformer. The weighted sum is
combined with a projection of the eight calibrated risks, which reach the
head whatever weight their agent received:
\begin{equation}
  u_t = \sum_{a} \alpha_{a,t} h_{a,t} + W_r r_t + b_r .
\end{equation}
A causal TCN with the same structure as the vitals agent (three residual
blocks, kernel $k=3$, dilations $\delta \in \{1, 2, 4\}$) then reads $u_1, \ldots, u_t$.
Each block adds $2(k-1)\delta$ hours of reach, so the head sees the current hour and
the 28 before it, a 29-hour receptive field. A linear layer gives the head's
logit. Heads are trained with binary cross-entropy over observed hours (padding
masked), Adam with learning rate $10^{-3}$, batches of 64 stays, at most 8
epochs and early stopping (patience 2) on validation AUPRC of the head's own
target. Training runs on one CPU thread so that runs are bit-identical. Each
head has 26{,}018 parameters.

\subsection{Two heads and noisy-OR}
A single head trained on the composite label must compromise between its
components: treatment labels carry information about timing, while death
labels carry most of the information needed to rank death. Sentinel therefore
trains two heads with identical architecture on the same agent bank. The
\emph{Treatment head} predicts ventilation or vasopressor onset within 24
hours; the \emph{Mortality head} predicts death within 24 hours. Each head's
raw score is calibrated with isotonic regression fitted on $V_{\text{fit}}$.
A small rank-preserving tie-break removes isotonic plateaus, which would
otherwise inflate alarm volume at a threshold. The two probabilities are
combined by noisy-OR,
\begin{equation}
  s_t = 1 - \bigl(1 - p^{T}_t\bigr)\bigl(1 - p^{M}_t\bigr),
\end{equation}
the probability that either event occurs if the two are independent. The
combination has no fitted weights \ev{B5}.

\subsection{Explanation layer}
When an alarm fires, a local language model (gemma4:12b, run through Ollama on
the same machine) writes one sentence explaining it. It receives the attention
weight and risk of each agent, the seven NEWS2 sub-scores and four-hour trends
of five vital signs. It is not given the risk score or told that an alarm
fired, and its output is never passed back to the model, so removing it leaves
every score unchanged \ev{J1}.

% ---------------------------------------------------------------------------
\section{Experimental setup}\label{sec:experiments}

\subsection{Baselines}
\paragraph{NEWS2.} Computed every hour from carried-forward vitals as above,
and thresholded with the same protocol as the models.

\paragraph{Equal-information GBM.} To separate the effect of the architecture
from the effect of data access, we train an XGBoost model on the same
information the agents read: vitals, laboratory values and their staleness,
plus 6-, 12- and 24-hour changes and trailing minima, maxima and means (145
features) \ev{B2}. A small grid over tree depth $\{4, 6, 8\}$ and class
weighting (none or fully rebalanced) is tuned on validation AUPRC, with
learning rate 0.05, up to 3{,}000 trees and early stopping after 50 rounds.

\paragraph{Published recipes.} We rebuild three published early-warning
models and train them on our cohort \ev{B3}. These are rebuilds of each
paper's feature recipe, not the commercial products.
\begin{itemize}
  \item \textbf{eCART-style GBM} \cite{churpek2025ecartv5}: current values,
        24-hour minima, maxima and changes of vitals, oxygen settings, GCS and
        laboratory values, plus time of day, age and sex (95 features);
        XGBoost trained on data down-sampled to 50\% positive hours, as in the
        paper.
  \item \textbf{YAIB LightGBM} \cite{vandewater2024yaib}: forward-filled
        signals with missingness indicators and running minimum, maximum,
        mean and count over the stay so far, plus age and sex (218 features);
        LightGBM over YAIB's published search space.
  \item \textbf{Sadanandan XGBoost} \cite{sadanandan2026multimodal}, the only
        published model we found with our exact composite label: current
        values of its vital and laboratory channels plus age and sex.
\end{itemize}
Each recipe's departures from the paper are listed in
Appendix~\ref{app:specs}. The most important are hourly scoring instead of
eCART's 8-hour survival blocks, and our in-ICU outcome instead of eCART's
ward outcome. All baselines are tuned on validation data, calibrated on
$V_{\text{fit}}$ and thresholded on $V_{\text{freeze}}$ exactly as Sentinel is.

\subsection{Seeds and inference}
Sentinel's headline results are means over ten training seeds of the fusion
heads (seeds 42 and 1--9), with the agent outputs held fixed. Baselines are
trained once. Head-to-head operating-point analyses use seed 42 unless stated
otherwise \ev{K4}. On eICU we additionally re-train every component (agents,
GBM and heads) under ten further seeds.

\subsection{External validation}
We evaluate on eICU-CRD v2.0 \todo{cite Pollard 2018}, a multi-centre US
database of 208 hospitals that are mostly community rather than academic
centres. The eICU cohort is built by a query that mirrors the MIMIC one and
reuses the same Python transformations, so the two datasets pass through one
code path \ev{G1}. Event definitions are mapped as closely as the data allow:
ventilation onset is the earlier of the recorded ventilator start and the
first mechanical-ventilation care-plan entry; vasopressor onset comes from
eICU's derived vasopressor table; and death is timed at ICU or hospital
discharge. Nothing is fitted on eICU: weights, normalisation statistics,
calibrators and thresholds all come from MIMIC. eICU has no order-entry data
from which to build the cadence agent, so the external model is the
seven-agent version of the two-head flagship, trained on MIMIC without
cadence. eICU also has no ECG data, so the ECG agent sends its no-data value
throughout. Because eICU records vitals less often than MIMIC, we report a
matched-density analysis restricted to hours in which all five vitals were
measured within the previous two hours, alongside the pooled result.
Per-hospital results are reported for the 91 hospitals with at least 20{,}000
scored hours and 50 positive hours.

We also measure the cost of adapting to a new site. For fine-tuning, eICU is
split by hospital into 129 training, 30 validation and 49 test hospitals, and
we compare recalibration only, retraining the two heads, and full fine-tuning
\ev{G7}. For calibration, we refit the final calibration on increasing numbers
of local patients \ev{G8}.

\subsection{Deployment measurements}
We time end-to-end scoring of one patient-hour (all agents and both heads,
with 29 hours of context) on one CPU thread on four laptops spanning three
operating systems and two processor architectures, with any discrete GPU
hidden or absent \ev{I2}. We also replay MIMIC test stays reading by reading,
to compare hourly scoring with scoring each new reading as it arrives
\ev{H1,H2}.

% ---------------------------------------------------------------------------
\section{Results}\label{sec:results}
\subsection{Discrimination}
% Table 1: AUPRC/AUROC with CIs for all baselines and the flagship
%   \ev{B1,B2,B3,B4}. Mortality head on death \ev{B5}.
\subsection{Patients caught at a fixed alarm budget}
% Figure: catch rate vs budget, 0.5--10\% \ev{C4} (source script must be
%   committed first). Cost: quiet-stay false alarms \ev{C3}.
\subsection{How agents should communicate}
% T20 LLM-text arm \ev{D1}. Same-fuser channel arms \ev{D2} -- PENDING RE-RUN
%   ON THE FLAGSHIP. This subsection's claim is not final.
\subsection{Ablations}
% Agent additions \ev{E1}; ECG variants \ev{E2}; specialists vs one model
%   \ev{E3}; per-unit thresholds \ev{E4}; YAIB TCN swap \ev{E6}.
\subsection{Timing}
% Lead vs NEWS2 \ev{F1}; death-stay analysis \ev{F2,F3}. State the modest size.
\subsection{External validation and adaptation}
% eICU zero-shot \ev{G1--G5}; fine-tuning \ev{G7}; calibration repair \ev{G8}.
\subsection{Real-time scoring}
% Reading-level replay \ev{H1}; vitals-only real time \ev{H2}.
\subsection{Deployment footprint and explanations}
% CPU latency across four machines \ev{I2,I3,I4}; card faithfulness \ev{J1,J2}.
\todo{draft}

% ---------------------------------------------------------------------------
\section{Discussion}\label{sec:discussion}
% What the evidence supports and what it does not; the validation ladder
% (retrospective single-site, retrospective external, prospective silent,
% independent replication).
\subsection{Limitations}
% \ev{K1--K7}: single development site, retrospective only; cadence reads
% clinician behaviour; no clinician review of per-modality design yet;
% single-seed experiments; hour-level metric; untested design choices (E9).
\todo{draft}

\section{Conclusion}\label{sec:conclusion}
\todo{draft}

% ---------------------------------------------------------------------------
\section*{Data and code availability}
% MIMIC-IV and eICU-CRD under the PhysioNet credentialed-access DUA; no
% patient-level data released. Code: \todo{repository link and commit}.

\bibliographystyle{plainnat}
% For arXiv: copy references.bib next to main.tex (or upload main.bbl).
\bibliography{../docs/references}

\appendix
\section{Cohort construction details}
\section{Agent inputs and hyperparameters}
\section{Per-seed results}
\section{Published-spec rebuilds: deviations from each paper}\label{app:specs}

\end{document}
